\chapter{Appendix 8: Putting the Web to Work Inside Your Business}
\label{web_work}

\section{Forword}

The Internet is one of the fastest growing computing transformations
ever. And although external, ``commercial'' web sites are the highest
profile applications on the Internet, an equally important and exciting
phenomenon is taking place behind corporate firewalls on internal,
corporate-wide networks. Internal web sites, which run on open TCP/IP
networks, are performing real business applications across
organizations by publishing information, enabling transactions, and
facilitating communities of interest. Across all industries,
institutions are beginning to use internal web applications to meet
departmental, interdepartmental and companywide communication needs and
improve user productivity.


\section{Introduction}

The growth of the World Wide Web has been explosive. Leading measures
estimate there are approximately 30 millions users on the Internet
today, with nearly 8-10 million of those users accessing the World Wide
Web. Web servers, which deliver the content to individuals (or clients)
have exploded from an estimated 130 Web servers as of June 1993 to over
15,700 web servers as of April, 1995 (Matthew Gray's WebCrawler Report).\\

The World Wide Web explosion is occurring because a critical mass of
enabling technologies and market conditions has met a pent-up demand
for a new way of linking people and information. The key enablers of
Web growth are:
\begin{itemize}
  \item The proliferation of PCs, LANs, and modems.
  \item Open standards such as TCP/IP, HTTP, and HTML.
  \item Cross-platform support.
  \item Multimedia support and ease of use.
  \item Support for secure transactions.
\end{itemize}

These market enablers touched off an explosion of demand for web
services, providing people with a new and better way of linking up with
information. Previously, the ``information age'' had inundated people
with data that was out of date, incomplete, or irrelevant. What
organizations need instead and what the web provides is a way for
people to easily access the information they need, when they need it.\\

External web sites offer individuals and corporations access to large
amounts of invaluable information. Many corporations have been quick to
realize, though, that these same benefits of easy and rich information
access can apply to the internal distribution of their corporate or
departmental information. Consequently, there has been a concurrent,
but less publicized, explosion of internal web sites and servers -
behind the corporate firewall.\\

Because internal web sites are just emerging, the current size and
growth rate of this market is difficult to determine. However, Mark
Winther, Internet analyst at IDC, estimates that ``in the short term -
over the next five years - the enterprise-side market for Internet
server technologies may grow slightly more slowly than the external,
commercial market for those technologies. However, it will more than
make up for this slow start: demand for enterprise applications of web
technology will probably end up being five times greater than for
external applications.''


\section{Internal Use of the Web: An Exploding Market}

Predictions for internal web growth spring from an examination of
current corporate needs and market demands. Organizations of all sizes
are facing major communication challenges, and in today's fast-moving
competitive environment, quick and easy access to departmental,
corporate, and customer information is essential.\\

Companies are realizing that many of their current communications
vehicles are too limited to cope with today's market environment.
Intuitive access to broad types of media make the web an ideal solution
to these communication problems. As more companies are realizing its
superiority, the number of internal web sites is quietly growing.
Enterprise web systems can provide information in a way that is:
\begin{itemize}
  \item immediate
  \item cost effective
  \item easy to use
  \item rich in format
  \item versatile
\end{itemize}

But the internal web does more than just solve the problems inherent in
existing communications vehicles. It provides:
\begin{itemize}
  \item Freedom of choice: web technology is based on open standards and
therefore doesn't lock companies into limited, costly choices. In fact,
web technology is available for nearly all leading operating systems
and hardware platforms and can also leverage legacy database systems.

  \item Security: protecting information, even within a private corporate
network, is critical. Netscape's secure socket layer (SSL) technology
encrypts packets of information so that information is transferred
securely. Netscape has openly licensed SSL, making it the leading de
facto security standard for web-related products.

  \item Ease of use: one of the most fundamental web technologies driving
more and more people to use web technology is hyperlinking.
Hyperlinking allows users to easily navigate and find information by
simply clicking on a word or graphic. In addition, with web clients
such as Netscape Navigator, a single front end is used to access all
internal and external resources - so users don't need to learn multiple
software packages. Ease of use is a critical factor for a corporate
information system - in fact, according to a recent survey, 56\% of
managers said this was the single most critical factor.

  \item Cost-effectiveness: web applications are surprisingly inexpensive,
both in initial purchase, training, and in deployment. At less than \$40
per user on an interdepartmental or companywide basis, they cost far
less than most other communications or workgroup systems.
\end{itemize}


\section{Improving Internal Communications}

Convinced by these advantages, aerospace, investment banking,
telecommunication, entertainment, manufacturing, computer software and
hardware companies to name a few have already implemented internal web
applications and report significant cost savings and improvements in
staff productivity. Although the range of applications that can be
developed to meet industry-specific or general needs is virtually
limitless, web applications generally fit in to one of three categories:
\begin{itemize}
  \item Publishing applications are essentially one-to-many
communications: teams, departments, or entire corporations can set up
pages where they post information, reducing bulky, easily outdated
paper-based information. Applications like this bring an immediate
payback to organizations, reducing the costs of producing, printing,
shipping, and updating corporate information.

  \item Transaction applications are two-way interactions, such as
downloading software or checking benefits information. Whether an
employee needs a report, a software download, or a personalized letter,
using web technology, linked to legacy data, can be an intuitive and
efficient alternative to the delays and frustrations of telephone tag
or paper pushing.

  \item Community applications are many-to-many interactions. They include
newsgroups that facilitate direct exchanges of information between
members of a group, making information available to others within the
group. Newsgroups aren't email aliases. They resemble a tabloid that
readers scan, drilling down to more data if they wish. People subscribe
to the newsgroups that interest them, and what they see on screen is a
full list of subject lines, authors, and news article numbers. Each of
these items is the beginning of a ``thread'', which starts when someone
sends out an article or e-mail.
\end{itemize}

These types of web applications can improve communications and
productivity across all areas of a company, although they may take very
different forms depending on whether they're designed to meet the needs
of specific departments or corporate-wide functions. Company
departments like sales and marketing, product developments and others
currently are deploying web applications to solve their communications
problems in the following ways.

\subsection{Sales and Marketing Applications}

One of the fundamental challenges in sales and marketing departments is
getting up-to-date and reference information to an often geographically
distributed group of people. Web publishing applications for sales and
marketing departments reflect the fact that having just the right
information at the right moment can make a sale happen, and lacking
that information can mean losing the deal to your competitor. Web
applications can be developed to give immediate access to the following
types of information:
\begin{itemize}
  \item Product information, pricing charts, and new collateral sales
leads.
  \item Key competitive information.
  \item Lists of key customer wins, including win/loss analysis.
  \item Program calendars so that sales is aware of planned marketing
activities and sales forecasts.
  \item On-line training materials.
\end{itemize}

Web transaction applications, coupled with industry standard databases,
can help marketing departments maintain and access a database of key
customers and references. Sales people can use transaction applications
to place orders, check order status, and complete the paperwork
surrounding a sale.\\

One Netscape customer in the pharmaceutical industry has an unusual
transaction application, which enables its sales representatives to
respond more quickly to customer questions. For example, during a sales
visit, a doctor may ask the representative for a letter from the
company explaining the risks of using one of its products during
pregnancy. Rather than spending hours trying to reach the appropriate
person at headquarters to get the letter sent, the representative can
call up a Netscape form, enter the doctor's name and address and the
information request, and submit that form to headquarters immediately.\\

In community types of applications, sales and marketing teams share
expertise that help them be more successful in specific sales
situations. Newsgroups can focus on:
\begin{itemize}
  \item Issues surrounding winning deals against a particular competitor:
weak points, rebuttals, and the latest tactics of the competitor as
experienced by sales people in the field.
  \item Feedback and discussion of various marketing programs their
effectiveness and ways in which they're working well and suggestions
for improvement.
  \item Threaded discussion on specific accounts or issues, to help sales
representatives to exchange information about the challenges of a
particular segment or customer.
\end{itemize}

\subsection{Product Development Applications}

Like sales and marketing departments, product development teams need
up-to-date information to do their jobs. Product development
applications often center on project management, with team members
updating project schedules and sharing information about the progress
of development or customer feedback. Of course, access to the sensitive
information in these applications is restricted to team members. The
types of information made accessible through web applications may
include:
\begin{itemize}
  \item Product specifications, product designs, product schedule
milestones and changes.
  \item Team member listings and responsibilities, key customer issues.
  \item Features of key competitive products, etc.
\end{itemize}

Transaction-type applications allow team members to request or submit
specific information. For example, a web-client and server application
used in conjunction with a database can be used to log and query
product test results. Team members could also use web applications to
submit a request to patent a recent development, sending the
appropriate forms and information directly to the corporate legal
department.\\

Community applications facilitate the discussions and exchange of ideas
that are such a critical part of product development, such as threads
focusing on:\\
Expertise on specific issues in product development.\\

A forum for brainstorming new ideas and critiquing proposed approaches.\\

Customer issues feedback to product development from sales, marketing,
and customer service.

\subsection{Customer Service and Support Applications}

The goal of customer service and support groups is to provide the best
quality service in the most cost-effective, efficient way possible.
Both publishing and transaction web applications enable customer
service and support departments to share customer and problem
information, creating a coordinated support system. Such applications
enable team members to:
\begin{itemize}
  \item Share up-to-date status reports on problems so that all team
members can respond to customer calls.
  \item Get the most up-to-date information on the status of customers'
orders.
  \item Be alerted immediately to any important changes, such as special
offers or problems.
  \item Train on-line to respond to customer queries and complaints.
\end{itemize}

Community applications enable service and support staff to delve more
deeply into specific issues. Threads provide a forum for detailed
discussion of underlying factors in a recurrent customer problem and
for an exchange of experience of what has succeeded or failed in
addressing a problem. In addition, newsgroups alert staff assigned to
particular accounts of changes or other important information.

\subsection{Human Resources Applications}

More than any other department within a company, Human Resource's job
is to keep employees well informed of important company issues in
addition to all the details pertaining to their personal health and
well being. By using web applications to publish corporate information
and using transaction type applications to provide personal data, human
resource departments can free their staffs from answering routine
questions and basic processing tasks. Human Resource departments can
publish information such as:
\begin{itemize}
  \item Corporate benefits information.
  \item Corporate policies.
  \item Company mission and goals.
  \item Job postings.
  \item Telephone directories.
  \item The annual report.
  \item Employee development.
\end{itemize}

Transaction applications eliminate much of the routine paper processing
that HR departments typically handle and provides employees access to
their personal information. Transactions, particularly secure
transactions can enable:
\begin{itemize}
  \item On-line employee enrollment into benefit plans, such as 401K
programs.
  \item Employee status checking on their vacation balances, options etc.
  \item Managers' submission of employee status change forms.
\end{itemize}

Community newsgroups hosted by HR departments are particularly valuable
within large companies, because they provide a quick and cost-effective
way of communicating with specific groups of people such as senior
managers, people with 401K accounts, and employees with dependents.
Such newsgroups simplify the task of letting these people know about
changes to corporate policies that would affect them and also provide a
forum for discussion among people with shared interests.

\subsection{Finance Applications}

Well-managed companies watch their financial indicators carefully and
set clear financial objectives for their managers. Providing access to
important financial information securely and in an easy-to-use, on-line
manner has always been a significant challenge. By using internal web
applications, Finance can begin to more easily disseminate this
information to key managers by securely ``posting'' corporate finance
information or by providing simple forms based query capabilities.\\

The purchasing side of financial operations can also benefit from web
applications. Using the same products which enable ``cybermalls'' on
the public Internet today, companies can simplify electronic software
distribution and billing and the purchasing of supplies by providing an
internal web-based mall which includes all approved products for
purchase.


\section{The Future Of Internal Web Applications}

Web-based internal applications are possible today and are quickly
evolving to become an important part of many leading companies' IS
solutions. A competitive and open software and hardware marketplace is
driving companies like Netscape to make improvements in the technology
to enhance the range of what it can do and how easily web applications
can be managed and deployed. As the Internet continues to mature into
the mainstream platform for computing, internal web sites (as well as
external ones) are requiring better application development tools,
middleware solutions for connecting to legacy systems, and system
management tools.\\

Better development tools aim to add sophistication and complexity to
types of internal web applications that companies build. Specifically,
tools are on the horizon that will enhance HTML authoring, so that
users can easily create and publish rich, ``live'' multimedia
documents. Others will help the system managers of the internal web
site manage links, expiration dates, and ownership: these site and
document management tools will become increasingly critical as web
sites grow. Even now, some companies have more than 250,000 documents
on their internal web. Improved application development tools will make
it easier to build applications on top of web platforms that conduct
complex transactions against databases, legacy systems, and other
corporate information sources.\\

Increasingly powerful management tools will also enable an enterprise
to support far-reaching and innovative systems without sacrificing
control. Improved security and administration tools will ensure that
users are authorized to access the information they seek. Information
management tools will also provide integrated search technology to
effectively search across distributed servers.\\

The platform for creating more complex customized applications is
addressed today by products such as Netscape Servers with the Netscape
Server API and the Netscape Navigator with Netscape Client API.
Netscape is developing innovative technology that will make it possible
to build different applications, supporting ``live'' on-line content.
This technology, based on the Java language, as well as the integration
of third-party multimedia formats such as Adobe Acrobat and Macromedia
Director, will make Netscape the platform for applications that provide
intelligent and rich content.


\section{Conclusion}

The web-based products provided by Netscape Communications as well as
other leading vendors are today being used to improve the way companies
link people and information both on the external, vastly popular World
Wide Web as well as within private companywide webs. Today's growth of
these products are driven by the open environment and powerful
capabilities they bring to solving the daily business problems faced by
companies in increasingly competitive environments. These powerful
technologies are being adopted by leading companies because they
deliver an unprecedented and powerful combination of openness and
security, intuitive access to highly graphical information, and
flexibility for customization.
